\documentclass[11pt,a4paper]{article}

\usepackage[utf8]{inputenc}
\usepackage[T1]{fontenc}
\usepackage{lmodern}
\usepackage[margin=2.3cm]{geometry}
\usepackage[version=4]{mhchem}
\usepackage{booktabs}
\usepackage{enumitem}
\usepackage[dvipsnames]{xcolor}
\usepackage{tcolorbox}
\usepackage{listings}
\usepackage{microtype}
\usepackage{qrcode}
\usepackage[colorlinks=true,linkcolor=MidnightBlue,urlcolor=MidnightBlue]{hyperref}

\lstset{basicstyle=\ttfamily\small,columns=fullflexible,keepspaces=true,
        frame=single,framerule=0.3pt,rulecolor=\color{gray!60},backgroundcolor=\color{gray!6}}

\newtcolorbox{keybox}[1]{colback=MidnightBlue!4,colframe=MidnightBlue!70,
  fonttitle=\bfseries,title=#1,boxrule=0.6pt,arc=2pt}
\newtcolorbox{warnbox}[1]{colback=BrickRed!4,colframe=BrickRed!70,
  fonttitle=\bfseries,title=#1,boxrule=0.6pt,arc=2pt}

% Colab menus and buttons
\newcommand{\menu}[1]{\textsf{\textbf{#1}}}

\title{\textbf{Running the MLIP tutorials on Google Colab}\\[4pt]
\large Student guide for the three notebooks of \texttt{TP-MLIP}}
\author{Machine-learned interatomic potentials --- hands-on session}
\date{}

\begin{document}
\maketitle

\begin{minipage}[c]{0.70\textwidth}
The three tutorials on machine-learned interatomic potentials (MLIPs) run on
\textbf{Google Colab}. All you need is a web browser and a Google account.
Nothing has to be installed on your computer: Colab gives you a Linux machine
with 2 CPUs, and the first cell of each notebook installs everything else on
it.

\medskip
Everything starts from the course page, which you can open with the QR code:
\begin{center}
\url{https://github.com/zacway/TP-MLIP}
\end{center}
\end{minipage}\hfill
\begin{minipage}[c]{0.25\textwidth}
\centering\qrcode[nolink,height=3.6cm]{https://github.com/zacway/TP-MLIP}
\end{minipage}
\medskip

%======================================================================
\section{The three tutorials}
%======================================================================

Do them in order. Tutorials~1 and~2 each come with lesson notes that explain
the theory and discuss the results: read them alongside the notebook. They are
linked from the course page.

\begin{center}
\begin{tabular}{@{}lll@{}}
\toprule
Tutorial & Notebook & Computing time\\
\midrule
1. Zundel ion, \ce{H5O2+}                   & \texttt{zundel.ipynb}     & $\sim$10 min\\
2. Protonated water cluster, \ce{(H2O)6H+}  & \texttt{wclusters.ipynb}  & $\sim$20 min\\
3. Published models for aqueous systems     & \texttt{wsystems.ipynb}   & $\sim$20 min\\
\bottomrule
\end{tabular}
\end{center}

The computing time is the total time spent waiting for cells to finish,
setup included. Most of it goes into training the neural networks and running
molecular dynamics (MD). Reading, thinking and answering the questions come on
top of it.

%======================================================================
\section{Step by step}
%======================================================================

\begin{enumerate}[leftmargin=*,itemsep=4pt]
  \item \textbf{Open the tutorial in Colab.} On the course page, click on the
        \menu{Open in Colab} button of the tutorial. The notebook opens in
        Colab; sign in with your Google account if asked.

  \item \textbf{Run the setup cell.} Click on the first code cell and press
        \menu{Shift+Enter}. Colab warns that the notebook was not authored
        by Google: click \menu{Run anyway}. The cell downloads the course files, installs n2p2 (to
        train the neural networks) and CP2K (to run MD), and moves into the
        tutorial's folder. It takes 1 to 2 minutes the first time, and prints
        \texttt{Setup complete.} when it is done.

        \begin{warnbox}{If the setup cell stops with ``Could not download TP-MLIP.zip''}
        Download \texttt{TP-MLIP.zip} yourself from the \menu{Releases}
        section of the course page (right-hand side). In Colab, click on the
        folder icon in the left sidebar to open the \menu{Files} panel, drag
        and drop the zip file into it, wait until the upload is finished, then
        run the setup cell again.
        \end{warnbox}

  \item \textbf{Work through the notebook from top to bottom.} Run the cells
        one by one with \menu{Shift+Enter}, reading the text as you go. The
        questions are written in blue. Each cell uses the results of the cells
        above it, so do not skip any.
        \begin{itemize}[nosep]
          \item MD cells call \texttt{run\_cp2k(...)}. While the simulation
                runs, they print its latest step every 15 seconds, then
                \texttt{Done in X min.} Wait for that line before moving on.
          \item Training cells take a few minutes as well. The
                active-learning cell prints a summary after each of its 15
                iterations; again, wait until the cell has finished.
        \end{itemize}

  \item \textbf{Tutorials 2 and 3.} Open each notebook the same way
        (step~1) and run its setup cell first. The tutorials do not depend on
        each other: each one trains or downloads its own models.
\end{enumerate}

%======================================================================
\section{Good to know}
%======================================================================

\paragraph{Keep the default CPU runtime.} Do not switch to a GPU in
\menu{Runtime > Change runtime type}: n2p2 and CP2K only use the CPU.

\paragraph{Looking at and editing files.} The tutorial files are in
\texttt{/content/TP-MLIP}, which you can browse in the \menu{Files} panel.
Double-click on a text file, such as a CP2K input (\texttt{*.inp}), to open it
in an editor next to the notebook, and press \menu{Ctrl+S} to save your
changes.

\paragraph{Colab's files are temporary.} After some time without activity
(typically around 90 minutes), when you close the tab for long, or after
\menu{Runtime > Disconnect and delete runtime}, Colab gives you a fresh
machine and everything computed so far is lost: trained models, trajectories,
installed programs. Your notebook and its displayed outputs remain. To carry on,
click on the cell where you stopped and use \menu{Runtime > Run before}: it
runs the setup cell and all the other cells above it again.

\paragraph{Saving your work.} The notebook opened from the course page is
not saved anywhere. \menu{File > Save a copy in Drive} keeps a copy of the
notebook with your notes and outputs (plots, scores), but not the files the
notebook produced. To keep one of those, use the \menu{$\vdots$} menu next to it
in the \menu{Files} panel and choose \menu{Download}.

\paragraph{Training again with other settings.} The notebooks invite you to
try other network architectures. Before training again, delete the previous
results, or they will interfere with the new ones. Run this in a new code cell
(\menu{+ Code}) before re-running the training cells:
\begin{lstlisting}
!rm -rf {base_dir}/qbc/iteration-* {base_dir}/qbc/results.shelf*
!rm -rf {base_dir}/train/final_training
\end{lstlisting}
and, for the single model at the start of tutorial~1:
\begin{lstlisting}
!rm -rf {base_dir}/single/training
\end{lstlisting}

\begin{keybox}{Changing the number of committee members}
The CP2K input files list the committee members one by one, in one
\texttt{\&MODEL} block each (8 blocks). If you change \texttt{n\_members}, edit
the \texttt{\&MODEL} blocks of the inputs used afterwards so that they match:
\texttt{nnp\_md/zundel.inp} and \texttt{cluster/cluster.inp} in tutorial~1;
\texttt{nnp\_md/cluster.inp}, \texttt{big\_cluster/big\_cluster.inp} and
\texttt{bulk/bulk.inp} in tutorial~2.
\end{keybox}

\paragraph{If something goes wrong.} Use \menu{Runtime > Restart session},
which keeps the files, then use \menu{Runtime > Run before} on the cell you
were at. If it is still broken, use \menu{Runtime > Disconnect
and delete runtime} and start again from the setup cell.

\end{document}
